% Indeks generowany na podstawie \begin{definicja|twierdzenie|lemat|wniosek}.
\chapter*{Indeks definicji i twierdzeń}
\markboth{INDEKS DEFINICJI I TWIERDZEŃ}{INDEKS DEFINICJI I TWIERDZEŃ}
\phantomsection\addcontentsline{toc}{chapter}{Indeks definicji i twierdzeń}
Indeks podaje nazwę rezultatu, jego numer i stronę. Nazwy i numery są aktywnymi odsyłaczami.
Uwzględniono także lematy i wnioski.

\section*{Definicje}
\begingroup\small
\setlength\LTleft{0pt}\setlength\LTright{0pt}
\begin{longtable}{@{}p{0.72\linewidth}p{0.135\linewidth}p{0.08\linewidth}@{}}
\textbf{Nazwa} & \textbf{Nr} & \textbf{Str.}\\\hline
\endfirsthead
\multicolumn{3}{@{}l}{\emph{Definicje -- ciąg dalszy}}\\[2pt]
\textbf{Nazwa} & \textbf{Nr} & \textbf{Str.}\\\hline
\endhead
\hyperref[def:atom]{Atom i miara bezatomowa} & \ref*{def:atom} & \pageref*{def:atom} \\
\hyperref[def:ac-measure-17]{Bezwzględna ciągłość miary} & \ref*{def:ac-measure-17} & \pageref*{def:ac-measure-17} \\
\hyperref[def:reduced-boundary-29]{Brzeg zredukowany} & \ref*{def:reduced-boundary-29} & \pageref*{def:reduced-boundary-29} \\
\hyperref[def:signed-integral-over-set-7]{Całka funkcji całkowalnej po zbiorze} & \ref*{def:signed-integral-over-set-7} & \pageref*{def:signed-integral-over-set-7} \\
\hyperref[def:integral-simple]{Całka funkcji prostej} & \ref*{def:integral-simple} & \pageref*{def:integral-simple} \\
\hyperref[def:complex-integral]{Całka funkcji zespolonej} & \ref*{def:complex-integral} & \pageref*{def:complex-integral} \\
\hyperref[def:integral-positive]{Całka Lebesgue'a funkcji nieujemnej} & \ref*{def:integral-positive} & \pageref*{def:integral-positive} \\
\hyperref[def:integral-over-set]{Całka po podzbiorze} & \ref*{def:integral-over-set} & \pageref*{def:integral-over-set} \\
\hyperref[def:tight-single-20]{Ciasność pojedynczej miary} & \ref*{def:tight-single-20} & \pageref*{def:tight-single-20} \\
\hyperref[def:cauchy-measure-9]{Ciąg Cauchy'ego według miary} & \ref*{def:cauchy-measure-9} & \pageref*{def:cauchy-measure-9} \\
\hyperref[def:positive-negative-parts]{Części dodatnia i ujemna} & \ref*{def:positive-negative-parts} & \pageref*{def:positive-negative-parts} \\
\hyperref[def:vitali-fine-24]{Drobne pokrycie Vitaliego} & \ref*{def:vitali-fine-24} & \pageref*{def:vitali-fine-24} \\
\hyperref[def:dual-lp-15]{Dualność i norma funkcjonału} & \ref*{def:dual-lp-15} & \pageref*{def:dual-lp-15} \\
\hyperref[def:HL-max-24]{Funkcja maksymalna Hardy'ego--Littlewooda} & \ref*{def:HL-max-24} & \pageref*{def:HL-max-24} \\
\hyperref[def:simple-function]{Funkcja prosta} & \ref*{def:simple-function} & \pageref*{def:simple-function} \\
\hyperref[idx:06:definicja:11]{Funkcja zespolona mierzalna} & \ref*{idx:06:definicja:11} & \pageref*{idx:06:definicja:11} \\
\hyperref[def:limsup-sets]{Granica górna i dolna zbiorów} & \ref*{def:limsup-sets} & \pageref*{def:limsup-sets} \\
\hyperref[def:topological-group-23]{Grupa topologiczna} & \ref*{def:topological-group-23} & \pageref*{def:topological-group-23} \\
\hyperref[def:normal-jacobian-28]{Jakobian normalny} & \ref*{def:normal-jacobian-28} & \pageref*{def:normal-jacobian-28} \\
\hyperref[def:volume-jacobian-27]{Jakobian objętościowy} & \ref*{def:volume-jacobian-27} & \pageref*{def:volume-jacobian-27} \\
\hyperref[def:k-jacobian-28]{Jakobian wymiaru $k$} & \ref*{def:k-jacobian-28} & \pageref*{def:k-jacobian-28} \\
\hyperref[def:uniform-integrability]{Jednostajna całkowalność} & \ref*{def:uniform-integrability} & \pageref*{def:uniform-integrability} \\
\hyperref[def:uniform-tight-22]{Jednostajna ciasność} & \ref*{def:uniform-tight-22} & \pageref*{def:uniform-tight-22} \\
\hyperref[def:monotone-class]{Klasa monotoniczna} & \ref*{def:monotone-class} & \pageref*{def:monotone-class} \\
\hyperref[def:measure-finiteness]{Klasy miar} & \ref*{def:measure-finiteness} & \pageref*{def:measure-finiteness} \\
\hyperref[def:hausdorff-measure-26]{Koszt pokrycia i miara Hausdorffa} & \ref*{def:hausdorff-measure-26} & \pageref*{def:hausdorff-measure-26} \\
\hyperref[def:linf-14]{Kres górny istotny i przestrzeń $L^\infty$} & \ref*{def:linf-14} & \pageref*{def:linf-14} \\
\hyperref[def:multiplicity-28]{Krotność odwzorowania} & \ref*{def:multiplicity-28} & \pageref*{def:multiplicity-28} \\
\hyperref[def:balls-24]{Kule i krotność pokrycia} & \ref*{def:balls-24} & \pageref*{def:balls-24} \\
\hyperref[def:haar-23]{Lewostronna i prawostronna miara Haara} & \ref*{def:haar-23} & \pageref*{def:haar-23} \\
\hyperref[def:local-majorant-11]{Lokalna dominacja całkowalna} & \ref*{def:local-majorant-11} & \pageref*{def:local-majorant-11} \\
\hyperref[def:max-measure-25]{Maksymalna funkcja miary} & \ref*{def:max-measure-25} & \pageref*{def:max-measure-25} \\
\hyperref[def:measure]{Miara} & \ref*{def:measure} & \pageref*{def:measure} \\
\hyperref[def:borel-locfinite-20]{Miara borelowska i lokalna skończoność} & \ref*{def:borel-locfinite-20} & \pageref*{def:borel-locfinite-20} \\
\hyperref[def:pushforward]{Miara obrazowa} & \ref*{def:pushforward} & \pageref*{def:pushforward} \\
\hyperref[def:radon-measure-20]{Miara Radona --- przyjęta konwencja} & \ref*{def:radon-measure-20} & \pageref*{def:radon-measure-20} \\
\hyperref[def:restriction]{Miara śladowa i ograniczenie} & \ref*{def:restriction} & \pageref*{def:restriction} \\
\hyperref[def:vector-measure-30]{Miara wektorowa} & \ref*{def:vector-measure-30} & \pageref*{def:vector-measure-30} \\
\hyperref[def:complex-measure-16]{Miara zespolona} & \ref*{def:complex-measure-16} & \pageref*{def:complex-measure-16} \\
\hyperref[def:outer-measure]{Miara zewnętrzna} & \ref*{def:outer-measure} & \pageref*{def:outer-measure} \\
\hyperref[def:lebesgue-outer]{Miara zewnętrzna Lebesgue'a} & \ref*{def:lebesgue-outer} & \pageref*{def:lebesgue-outer} \\
\hyperref[def:metric-outer-21]{Miara zewnętrzna metryczna} & \ref*{def:metric-outer-21} & \pageref*{def:metric-outer-21} \\
\hyperref[def:complete-measure]{Miara zupełna} & \ref*{def:complete-measure} & \pageref*{def:complete-measure} \\
\hyperref[def:extended-measurable]{Mierzalność funkcji rozszerzonej} & \ref*{def:extended-measurable} & \pageref*{def:extended-measurable} \\
\hyperref[def:normal-BV-29]{Normalna z pochodnej miary} & \ref*{def:normal-BV-29} & \pageref*{def:normal-BV-29} \\
\hyperref[def:support-radon-20]{Nośnik miary} & \ref*{def:support-radon-20} & \pageref*{def:support-radon-20} \\
\hyperref[def:finite-perimeter-29]{Obwód w sensie De Giorgiego} & \ref*{def:finite-perimeter-29} & \pageref*{def:finite-perimeter-29} \\
\hyperref[def:measurable-map]{Odwzorowanie mierzalne} & \ref*{def:measurable-map} & \pageref*{def:measurable-map} \\
\hyperref[def:algebry]{Pierścień, algebra, $\sigma$-algebra} & \ref*{def:algebry} & \pageref*{def:algebry} \\
\hyperref[def:five-convergences]{Pięć rodzajów zbieżności} & \ref*{def:five-convergences} & \pageref*{def:five-convergences} \\
\hyperref[def:RN-derivative-18]{Pochodna Radona--Nikodyma} & \ref*{def:RN-derivative-18} & \pageref*{def:RN-derivative-18} \\
\hyperref[def:current-29]{Prąd i jego brzeg} & \ref*{def:current-29} & \pageref*{def:current-29} \\
\hyperref[def:premeasure]{Premiara} & \ref*{def:premeasure} & \pageref*{def:premeasure} \\
\hyperref[def:halfopen-box]{Prostopadłościan półotwarty} & \ref*{def:halfopen-box} & \pageref*{def:halfopen-box} \\
\hyperref[def:complex-measures-space-16]{Przestrzeń $\mathcal M(X,\mathcal F)$} & \ref*{def:complex-measures-space-16} & \pageref*{def:complex-measures-space-16} \\
\hyperref[def:standard-borel]{Przestrzeń polska i standardowa borelowska} & \ref*{def:standard-borel} & \pageref*{def:standard-borel} \\
\hyperref[def:measure-boundary-29]{Punkty gęstości i brzeg miarowy} & \ref*{def:measure-boundary-29} & \pageref*{def:measure-boundary-29} \\
\hyperref[def:weak-lp-15]{Refleksywność i zbieżność słaba} & \ref*{def:weak-lp-15} & \pageref*{def:weak-lp-15} \\
\hyperref[def:regularity-20]{Regularność względem topologii} & \ref*{def:regularity-20} & \pageref*{def:regularity-20} \\
\hyperref[def:ae-functions]{Równość i zbieżność prawie wszędzie} & \ref*{def:ae-functions} & \pageref*{def:ae-functions} \\
\hyperref[def:lp-14]{Równość prawie wszędzie i przestrzeń $L^p$} & \ref*{def:lp-14} & \pageref*{def:lp-14} \\
\hyperref[def:frechet-rademacher-28]{Różniczkowalność Frécheta} & \ref*{def:frechet-rademacher-28} & \pageref*{def:frechet-rademacher-28} \\
\hyperref[def:borel]{$\sigma$-algebra borelowska} & \ref*{def:borel} & \pageref*{def:borel} \\
\hyperref[def:sigma-generated]{$\sigma$-algebra generowana} & \ref*{def:sigma-generated} & \pageref*{def:sigma-generated} \\
\hyperref[def:sigma-of-function]{$\sigma$-algebra generowana przez odwzorowanie} & \ref*{def:sigma-of-function} & \pageref*{def:sigma-of-function} \\
\hyperref[def:product-sigma]{$\sigma$-algebra produktowa} & \ref*{def:product-sigma} & \pageref*{def:product-sigma} \\
\hyperref[idx:01:definicja:04]{$\sigma$-algebra śladowa} & \ref*{idx:01:definicja:04} & \pageref*{idx:01:definicja:04} \\
\hyperref[idx:03:definicja:09]{\texorpdfstring{$\sigma$}{sigma}-skończona premiara} & \ref*{idx:03:definicja:09} & \pageref*{idx:03:definicja:09} \\
\hyperref[def:bochner-30]{Silna mierzalność i całka Bochnera} & \ref*{def:bochner-30} & \pageref*{def:bochner-30} \\
\hyperref[def:signed-measure-16]{Skończona miara ze znakiem} & \ref*{def:signed-measure-16} & \pageref*{def:signed-measure-16} \\
\hyperref[idx:21:definicja:08]{Skończona zespolona miara Radona} & \ref*{idx:21:definicja:08} & \pageref*{idx:21:definicja:08} \\
\hyperref[def:cantor-set]{Środkowotrzeci zbiór Cantora} & \ref*{def:cantor-set} & \pageref*{def:cantor-set} \\
\hyperref[def:test-19]{Test losowy i jego moc} & \ref*{def:test-19} & \pageref*{def:test-19} \\
\hyperref[def:energy-26]{$t$-energia miary} & \ref*{def:energy-26} & \pageref*{def:energy-26} \\
\hyperref[def:pi-lambda]{Układ $\pi$ i układ Dynkina} & \ref*{def:pi-lambda} & \pageref*{def:pi-lambda} \\
\hyperref[def:variation-signed-16]{Wariacja całkowita} & \ref*{def:variation-signed-16} & \pageref*{def:variation-signed-16} \\
\hyperref[def:BV-variation-29]{Wariacja funkcji} & \ref*{def:BV-variation-29} & \pageref*{def:BV-variation-29} \\
\hyperref[def:vector-variation-30]{Wariacja i półwariacja} & \ref*{def:vector-variation-30} & \pageref*{def:vector-variation-30} \\
\hyperref[def:lebesgue-point-25]{Wartość średnia i punkt Lebesgue'a} & \ref*{def:lebesgue-point-25} & \pageref*{def:lebesgue-point-25} \\
\hyperref[def:conditional-exp-18]{Warunkowa wartość oczekiwana} & \ref*{def:conditional-exp-18} & \pageref*{def:conditional-exp-18} \\
\hyperref[def:almost-everywhere]{Własność prawie wszędzie} & \ref*{def:almost-everywhere} & \pageref*{def:almost-everywhere} \\
\hyperref[def:RNP-30]{Własność Radona--Nikodyma} & \ref*{def:RNP-30} & \pageref*{def:RNP-30} \\
\hyperref[def:hausdorff-dimension-26]{Wymiar Hausdorffa} & \ref*{def:hausdorff-dimension-26} & \pageref*{def:hausdorff-dimension-26} \\
\hyperref[def:singular-16]{Wzajemna osobliwość} & \ref*{def:singular-16} & \pageref*{def:singular-16} \\
\hyperref[def:equivalent-measures-17]{Wzajemna równoważność miar} & \ref*{def:equivalent-measures-17} & \pageref*{def:equivalent-measures-17} \\
\hyperref[def:vague-weakstar-22]{Zbieżność słaba-* i mglista} & \ref*{def:vague-weakstar-22} & \pageref*{def:vague-weakstar-22} \\
\hyperref[def:narrow-measures-22]{Zbieżność wąska} & \ref*{def:narrow-measures-22} & \pageref*{def:narrow-measures-22} \\
\hyperref[def:convergence-measure]{Zbieżność według miary} & \ref*{def:convergence-measure} & \pageref*{def:convergence-measure} \\
\hyperref[def:lp-convergence-9]{Zbieżność w średniej rzędu $p$} & \ref*{def:lp-convergence-9} & \pageref*{def:lp-convergence-9} \\
\hyperref[def:rectifiable-29]{Zbiory przeliczalnie prostowalne} & \ref*{def:rectifiable-29} & \pageref*{def:rectifiable-29} \\
\hyperref[def:continuity-set-22]{Zbiór ciągłości miary} & \ref*{def:continuity-set-22} & \pageref*{def:continuity-set-22} \\
\hyperref[def:positive-negative-16]{Zbiór dodatni i ujemny} & \ref*{def:positive-negative-16} & \pageref*{def:positive-negative-16} \\
\hyperref[def:caratheodory-measurable]{Zbiór mierzalny w sensie Carathéodory'ego} & \ref*{def:caratheodory-measurable} & \pageref*{def:caratheodory-measurable} \\
\end{longtable}
\endgroup

\section*{Twierdzenia}
\begingroup\small
\setlength\LTleft{0pt}\setlength\LTright{0pt}
\begin{longtable}{@{}p{0.72\linewidth}p{0.135\linewidth}p{0.08\linewidth}@{}}
\textbf{Nazwa} & \textbf{Nr} & \textbf{Str.}\\\hline
\endfirsthead
\multicolumn{3}{@{}l}{\emph{Twierdzenia -- ciąg dalszy}}\\[2pt]
\textbf{Nazwa} & \textbf{Nr} & \textbf{Str.}\\\hline
\endhead
\hyperref[thm:positive-additivity]{Addytywność całki nieujemnej} & \ref*{thm:positive-additivity} & \pageref*{thm:positive-additivity} \\
\hyperref[thm:cc-approx-in-measure-10]{Aproksymacja funkcji o nośniku mierzalnym skończonej miary} & \ref*{thm:cc-approx-in-measure-10} & \pageref*{thm:cc-approx-in-measure-10} \\
\hyperref[thm:simple-real-approx]{Aproksymacja funkcji o wartościach rzeczywistych} & \ref*{thm:simple-real-approx} & \pageref*{thm:simple-real-approx} \\
\hyperref[thm:lebesgue-outer-regular]{Aproksymacja otwartymi zbiorami} & \ref*{thm:lebesgue-outer-regular} & \pageref*{thm:lebesgue-outer-regular} \\
\hyperref[thm:lebesgue-inner-regular]{Aproksymacja zbiorami zwartymi} & \ref*{thm:lebesgue-inner-regular} & \pageref*{thm:lebesgue-inner-regular} \\
\hyperref[thm:finite-metric-regular-20]{Aproksymacja zbiorów borelowskich} & \ref*{thm:finite-metric-regular-20} & \pageref*{thm:finite-metric-regular-20} \\
\hyperref[thm:atom-sigma-finite]{Atomy miar $\sigma$-skończonych mają skończoną masę} & \ref*{thm:atom-sigma-finite} & \pageref*{thm:atom-sigma-finite} \\
\hyperref[thm:bayes-finite-19]{Bayes dla skończonej przestrzeni parametrów} & \ref*{thm:bayes-finite-19} & \pageref*{thm:bayes-finite-19} \\
\hyperref[thm:image-ac-19]{Bezwzględna ciągłość zachowana przez obraz} & \ref*{thm:image-ac-19} & \pageref*{thm:image-ac-19} \\
\hyperref[thm:lebesgue-completion]{Borelowskie aproksymacje i uzupełnienie} & \ref*{thm:lebesgue-completion} & \pageref*{thm:lebesgue-completion} \\
\hyperref[thm:borel-Rd]{Borelowskie prostokąty} & \ref*{thm:borel-Rd} & \pageref*{thm:borel-Rd} \\
\hyperref[thm:moving-domains-11]{Całkowanie po zależnej od parametru dziedzinie} & \ref*{thm:moving-domains-11} & \pageref*{thm:moving-domains-11} \\
\hyperref[thm:change-measure-19]{Całkowanie po zmianie miary} & \ref*{thm:change-measure-19} & \pageref*{thm:change-measure-19} \\
\hyperref[thm:sum-absolutely-integrable]{Całkowanie szeregów absolutnie sumowalnych} & \ref*{thm:sum-absolutely-integrable} & \pageref*{thm:sum-absolutely-integrable} \\
\hyperref[thm:pushforward-integral-13]{Całkowanie względem miary obrazowej} & \ref*{thm:pushforward-integral-13} & \pageref*{thm:pushforward-integral-13} \\
\hyperref[thm:product-projections]{Charakteryzacja przez współrzędne} & \ref*{thm:product-projections} & \pageref*{thm:product-projections} \\
\hyperref[thm:polish-tight-20]{Ciasność skończonych miar na przestrzeniach polskich} & \ref*{thm:polish-tight-20} & \pageref*{thm:polish-tight-20} \\
\hyperref[thm:parameter-continuity]{Ciągłość pod znakiem całki} & \ref*{thm:parameter-continuity} & \pageref*{thm:parameter-continuity} \\
\hyperref[thm:translation-lp-14]{Ciągłość translacji w $L^p(\mathbb R^d)$} & \ref*{thm:translation-lp-14} & \pageref*{thm:translation-lp-14} \\
\hyperref[thm:measure-continuity]{Ciągłość z dołu i z góry} & \ref*{thm:measure-continuity} & \pageref*{thm:measure-continuity} \\
\hyperref[thm:atomic-cont-singular-17]{Dalszy rozkład miary osobliwej na prostej} & \ref*{thm:atomic-cont-singular-17} & \pageref*{thm:atomic-cont-singular-17} \\
\hyperref[thm:weak-lsc-15]{Dolna półciągłość normy} & \ref*{thm:weak-lsc-15} & \pageref*{thm:weak-lsc-15} \\
\hyperref[thm:perimeter-lsc-29]{Dolna półciągłość obwodu} & \ref*{thm:perimeter-lsc-29} & \pageref*{thm:perimeter-lsc-29} \\
\hyperref[thm:dual-lp-15]{Dualność $L^p$ dla $1<p<\infty$} & \ref*{thm:dual-lp-15} & \pageref*{thm:dual-lp-15} \\
\hyperref[thm:dual-l1-15]{Dualność $L^1$ przy $\sigma$-skończoności} & \ref*{thm:dual-l1-15} & \pageref*{thm:dual-l1-15} \\
\hyperref[thm:tv-dual-tests-16]{Dualny opis normy w wariacji całkowitej} & \ref*{thm:tv-dual-tests-16} & \pageref*{thm:tv-dual-tests-16} \\
\hyperref[thm:ui-two-criteria-8]{Dwa równoważne opisy jednostajnej całkowalności} & \ref*{thm:ui-two-criteria-8} & \pageref*{thm:ui-two-criteria-8} \\
\hyperref[thm:measurable-algebra]{Działania algebraiczne} & \ref*{thm:measurable-algebra} & \pageref*{thm:measurable-algebra} \\
\hyperref[thm:total-variation-formula-19]{Formuła wariacyjna dla dwóch prawdopodobieństw} & \ref*{thm:total-variation-formula-19} & \pageref*{thm:total-variation-formula-19} \\
\hyperref[thm:frostman-26]{Frostmana, wersja dla zbiorów zwartych} & \ref*{thm:frostman-26} & \pageref*{thm:frostman-26} \\
\hyperref[thm:fubini-12]{Fubini} & \ref*{thm:fubini-12} & \pageref*{thm:fubini-12} \\
\hyperref[thm:cdf-derivative-25]{Funkcja rozkładu a absolutna ciągłość miary} & \ref*{thm:cdf-derivative-25} & \pageref*{thm:cdf-derivative-25} \\
\hyperref[thm:measurable-real]{Funkcje rzeczywiste i działania} & \ref*{thm:measurable-real} & \pageref*{thm:measurable-real} \\
\hyperref[thm:pushforward-composition]{Funktorialność miar obrazowych} & \ref*{thm:pushforward-composition} & \pageref*{thm:pushforward-composition} \\
\hyperref[thm:cc-dense-lp-radon-20]{Gęstość $C_c(X)$ w przestrzeniach $L^p$} & \ref*{thm:cc-dense-lp-radon-20} & \pageref*{thm:cc-dense-lp-radon-20} \\
\hyperref[thm:cc-in-lp-10]{Gęstość funkcji gładkich} & \ref*{thm:cc-in-lp-10} & \pageref*{thm:cc-in-lp-10} \\
\hyperref[thm:cc-dense-lp-14]{Gęstość funkcji gładkich w przestrzeni euklidesowej} & \ref*{thm:cc-dense-lp-14} & \pageref*{thm:cc-dense-lp-14} \\
\hyperref[thm:simple-dense-lp-14]{Gęstość funkcji prostych} & \ref*{thm:simple-dense-lp-14} & \pageref*{thm:simple-dense-lp-14} \\
\hyperref[thm:pushforward-coarea-28]{Gęstość obrazu przez submersję} & \ref*{thm:pushforward-coarea-28} & \pageref*{thm:pushforward-coarea-28} \\
\hyperref[thm:smooth-perimeter-29]{Granica klasycznej geometrii} & \ref*{thm:smooth-perimeter-29} & \pageref*{thm:smooth-perimeter-29} \\
\hyperref[thm:haar-main-23]{Haara: istnienie i jednoznaczność} & \ref*{thm:haar-main-23} & \pageref*{thm:haar-main-23} \\
\hyperref[thm:laplace-holo-11]{Holomorficzność transformaty Laplace'a} & \ref*{thm:laplace-holo-11} & \pageref*{thm:laplace-holo-11} \\
\hyperref[thm:holder-14]{Hölder} & \ref*{thm:holder-14} & \pageref*{thm:holder-14} \\
\hyperref[thm:finiteness-implications]{Implikacje} & \ref*{thm:finiteness-implications} & \pageref*{thm:finiteness-implications} \\
\hyperref[thm:lebesgue-exists]{Istnienie i podstawowe własności} & \ref*{thm:lebesgue-exists} & \pageref*{thm:lebesgue-exists} \\
\hyperref[thm:vitali]{Istnienie zbioru niemierzalnego} & \ref*{thm:vitali} & \pageref*{thm:vitali} \\
\hyperref[thm:CE-existence-18]{Istnienie, jednoznaczność i własności podstawowe} & \ref*{thm:CE-existence-18} & \pageref*{thm:CE-existence-18} \\
\hyperref[thm:BV-layercake-29]{Jedna nierówność formuły koobszaru dla $BV$} & \ref*{thm:BV-layercake-29} & \pageref*{thm:BV-layercake-29} \\
\hyperref[thm:uniform-convexity-15]{Jednostajna wypukłość} & \ref*{thm:uniform-convexity-15} & \pageref*{thm:uniform-convexity-15} \\
\hyperref[thm:haar-unique-23]{Jednoznaczność miary Haara} & \ref*{thm:haar-unique-23} & \pageref*{thm:haar-unique-23} \\
\hyperref[thm:uniqueness-cc-20]{Jednoznaczność miary Radona przez całki} & \ref*{thm:uniqueness-cc-20} & \pageref*{thm:uniqueness-cc-20} \\
\hyperref[thm:uniqueness-finite]{Jednoznaczność miar skończonych} & \ref*{thm:uniqueness-finite} & \pageref*{thm:uniqueness-finite} \\
\hyperref[thm:uniqueness-sigma-finite]{Jednoznaczność przy wspólnym wyczerpaniu} & \ref*{thm:uniqueness-sigma-finite} & \pageref*{thm:uniqueness-sigma-finite} \\
\hyperref[thm:extension-uniqueness]{Jednoznaczność rozszerzenia} & \ref*{thm:extension-uniqueness} & \pageref*{thm:extension-uniqueness} \\
\hyperref[thm:egorov-9]{Jegorow} & \ref*{thm:egorov-9} & \pageref*{thm:egorov-9} \\
\hyperref[thm:jensen-14]{Jensen dla miary probabilistycznej} & \ref*{thm:jensen-14} & \pageref*{thm:jensen-14} \\
\hyperref[thm:hausdorff-calibration-26]{Kalibracja wymiaru całkowitego bez twierdzeń geometrycznych} & \ref*{thm:hausdorff-calibration-26} & \pageref*{thm:hausdorff-calibration-26} \\
\hyperref[thm:bochner-basic-30]{Konstrukcja i oszacowanie całki Bochnera} & \ref*{thm:bochner-basic-30} & \pageref*{thm:bochner-basic-30} \\
\hyperref[thm:product-measure-12]{Konstrukcja miary produktowej} & \ref*{thm:product-measure-12} & \pageref*{thm:product-measure-12} \\
\hyperref[thm:outer-from-cover]{Konstrukcja przez pokrycia} & \ref*{thm:outer-from-cover} & \pageref*{thm:outer-from-cover} \\
\hyperref[thm:local-to-global-measure-9]{Kontrola ucieczki masy przy zbieżności lokalnej} & \ref*{thm:local-to-global-measure-9} & \pageref*{thm:local-to-global-measure-9} \\
\hyperref[thm:pointwise-limits]{Kresy i granice} & \ref*{thm:pointwise-limits} & \pageref*{thm:pointwise-limits} \\
\hyperref[thm:integrability-criterion]{Kryterium absolutnej całkowalności} & \ref*{thm:integrability-criterion} & \pageref*{thm:integrability-criterion} \\
\hyperref[thm:premeasure-continuity]{Kryterium ciągłości premiary skończonej} & \ref*{thm:premeasure-continuity} & \pageref*{thm:premeasure-continuity} \\
\hyperref[thm:eps-delta-ac-17]{Kryterium epsilon--delta} & \ref*{thm:eps-delta-ac-17} & \pageref*{thm:eps-delta-ac-17} \\
\hyperref[thm:tails-11]{Kryterium jednostajnej kontroli ogonów} & \ref*{thm:tails-11} & \pageref*{thm:tails-11} \\
\hyperref[thm:riemann-criterion]{Kryterium Lebesgue'a całkowalności Riemanna} & \ref*{thm:riemann-criterion} & \pageref*{thm:riemann-criterion} \\
\hyperref[thm:measurable-generator]{Kryterium przez generatory} & \ref*{thm:measurable-generator} & \pageref*{thm:measurable-generator} \\
\hyperref[thm:weak-lp-15]{Kryterium słabej zbieżności w $L^p$} & \ref*{thm:weak-lp-15} & \pageref*{thm:weak-lp-15} \\
\hyperref[thm:param-vitali-11]{Kryterium Vitaliego dla parametru} & \ref*{thm:param-vitali-11} & \pageref*{thm:param-vitali-11} \\
\hyperref[thm:series-lp-14]{Kryterium zbieżności szeregów w $L^p$} & \ref*{thm:series-lp-14} & \pageref*{thm:series-lp-14} \\
\hyperref[thm:lebesgue-diff-25]{Lebesgue'a o różniczkowaniu} & \ref*{thm:lebesgue-diff-25} & \pageref*{thm:lebesgue-diff-25} \\
\hyperref[thm:borel-cantelli-1]{Lemat Borela--Cantellego, część pierwsza} & \ref*{thm:borel-cantelli-1} & \pageref*{thm:borel-cantelli-1} \\
\hyperref[thm:fatou]{Lemat Fatou} & \ref*{thm:fatou} & \pageref*{thm:fatou} \\
\hyperref[thm:fatou-integrable-lower-8]{Lemat Fatou z całkowalnym ograniczeniem od dołu} & \ref*{thm:fatou-integrable-lower-8} & \pageref*{thm:fatou-integrable-lower-8} \\
\hyperref[thm:neyman-pearson-19]{Lemat Neymana--Pearsona} & \ref*{thm:neyman-pearson-19} & \pageref*{thm:neyman-pearson-19} \\
\hyperref[thm:5r-24]{Lemat o pięciokrotnym powiększeniu} & \ref*{thm:5r-24} & \pageref*{thm:5r-24} \\
\hyperref[thm:scheffe-9]{Lemat Scheffégo} & \ref*{thm:scheffe-9} & \pageref*{thm:scheffe-9} \\
\hyperref[thm:integral-linear-real]{Liniowość całki funkcji całkowalnych} & \ref*{thm:integral-linear-real} & \pageref*{thm:integral-linear-real} \\
\hyperref[thm:integral-complex-linear]{Liniowość zespolona i oszacowanie modułu} & \ref*{thm:integral-complex-linear} & \pageref*{thm:integral-complex-linear} \\
\hyperref[thm:linear-change-13]{Liniowy wzór na objętość} & \ref*{thm:linear-change-13} & \pageref*{thm:linear-change-13} \\
\hyperref[thm:local-convergence-measure]{Lokalna zbieżność według miary} & \ref*{thm:local-convergence-measure} & \pageref*{thm:local-convergence-measure} \\
\hyperref[thm:local-continuity-11]{Lokalne kryterium ciągłości} & \ref*{thm:local-continuity-11} & \pageref*{thm:local-continuity-11} \\
\hyperref[thm:null-local-27]{Lokalne kryterium zbiorów zerowych} & \ref*{thm:null-local-27} & \pageref*{thm:null-local-27} \\
\hyperref[thm:lch-secondcountable-radon-20]{Lokalnie skończone miary borelowskie} & \ref*{thm:lch-secondcountable-radon-20} & \pageref*{thm:lch-secondcountable-radon-20} \\
\hyperref[thm:lusin-10]{Łuzin} & \ref*{thm:lusin-10} & \pageref*{thm:lusin-10} \\
\hyperref[thm:measure-metric-9]{Metryka zbieżności według miary} & \ref*{thm:measure-metric-9} & \pageref*{thm:measure-metric-9} \\
\hyperref[thm:graph-area-28]{Miara na gładkim wykresie} & \ref*{thm:graph-area-28} & \pageref*{thm:graph-area-28} \\
\hyperref[thm:riemann-radon-27]{Miara objętościowa jest miarą Radona} & \ref*{thm:riemann-radon-27} & \pageref*{thm:riemann-radon-27} \\
\hyperref[thm:product-riemann-27]{Miara produktu riemannowskiego} & \ref*{thm:product-riemann-27} & \pageref*{thm:product-riemann-27} \\
\hyperref[thm:rn-manifold-27]{Miary absolutnie ciągłe względem objętości} & \ref*{thm:rn-manifold-27} & \pageref*{thm:rn-manifold-27} \\
\hyperref[thm:restriction-meas-5]{Mierzalność ograniczenia} & \ref*{thm:restriction-meas-5} & \pageref*{thm:restriction-meas-5} \\
\hyperref[thm:minimize-lp-15]{Minimalizacja wypukłej energii} & \ref*{thm:minimize-lp-15} & \pageref*{thm:minimize-lp-15} \\
\hyperref[thm:initial-sigma]{Minimalna struktura mierzalna} & \ref*{thm:initial-sigma} & \pageref*{thm:initial-sigma} \\
\hyperref[thm:minkowski-14]{Minkowski} & \ref*{thm:minkowski-14} & \pageref*{thm:minkowski-14} \\
\hyperref[thm:null-modifications]{Modyfikacje na zbiorach zerowych} & \ref*{thm:null-modifications} & \pageref*{thm:null-modifications} \\
\hyperref[thm:simple-up]{Monotoniczna aproksymacja funkcji nieujemnej} & \ref*{thm:simple-up} & \pageref*{thm:simple-up} \\
\hyperref[thm:fatou-for-sets]{Nierówności dla granic zbiorów} & \ref*{thm:fatou-for-sets} & \pageref*{thm:fatou-for-sets} \\
\hyperref[thm:lp-implies-measure-9]{Nierówność Czebyszewa i zbieżność według miary} & \ref*{thm:lp-implies-measure-9} & \pageref*{thm:lp-implies-measure-9} \\
\hyperref[thm:lp-interpolation-14]{Nierówność interpolacyjna} & \ref*{thm:lp-interpolation-14} & \pageref*{thm:lp-interpolation-14} \\
\hyperref[thm:markov-nonnegative]{Nierówność Markowa dla całki} & \ref*{thm:markov-nonnegative} & \pageref*{thm:markov-nonnegative} \\
\hyperref[thm:HL-weak-24]{Nierówność słabego typu $(1,1)$} & \ref*{thm:HL-weak-24} & \pageref*{thm:HL-weak-24} \\
\hyperref[thm:gibbs-19]{Nieujemność entropii względnej} & \ref*{thm:gibbs-19} & \pageref*{thm:gibbs-19} \\
\hyperref[thm:ae-integral-invariance]{Niezmienniczość całki przy zmianie na zbiorze zerowym} & \ref*{thm:ae-integral-invariance} & \pageref*{thm:ae-integral-invariance} \\
\hyperref[thm:support-fullmass-20]{Nośnik na przestrzeni drugiej przeliczalności} & \ref*{thm:support-fullmass-20} & \pageref*{thm:support-fullmass-20} \\
\hyperref[thm:riesz-c0-positive-21]{Ograniczone dodatnie funkcjonały na $C_0$} & \ref*{thm:riesz-c0-positive-21} & \pageref*{thm:riesz-c0-positive-21} \\
\hyperref[thm:extended-operations]{Operacje na funkcjach rozszerzonych} & \ref*{thm:extended-operations} & \pageref*{thm:extended-operations} \\
\hyperref[thm:partition-variation-16]{Opis wariacji przez podziały} & \ref*{thm:partition-variation-16} & \pageref*{thm:partition-variation-16} \\
\hyperref[thm:bounded-integral-complex-16]{Oszacowanie całki względem miary zespolonej} & \ref*{thm:bounded-integral-complex-16} & \pageref*{thm:bounded-integral-complex-16} \\
\hyperref[thm:frostman-energy-26]{Oszacowanie energii przez warunek Frostmana} & \ref*{thm:frostman-energy-26} & \pageref*{thm:frostman-energy-26} \\
\hyperref[thm:weak-measure-25]{Oszacowanie słabego typu dla miar} & \ref*{thm:weak-measure-25} & \pageref*{thm:weak-measure-25} \\
\hyperref[thm:mct]{O zbieżności monotonicznej; Beppo Levi} & \ref*{thm:mct} & \pageref*{thm:mct} \\
\hyperref[thm:dct]{O zbieżności zmajoryzowanej; Lebesgue} & \ref*{thm:dct} & \pageref*{thm:dct} \\
\hyperref[thm:BV-coarea-full-29]{Pełna formuła koobszaru dla $BV$} & \ref*{thm:BV-coarea-full-29} & \pageref*{thm:BV-coarea-full-29} \\
\hyperref[thm:local-differentiation-11]{Pierwsza pochodna — wersja lokalna} & \ref*{thm:local-differentiation-11} & \pageref*{thm:local-differentiation-11} \\
\hyperref[thm:nonnegative-first-properties]{Pierwsze własności całki nieujemnej} & \ref*{thm:nonnegative-first-properties} & \pageref*{thm:nonnegative-first-properties} \\
\hyperref[thm:BV-measure-29]{Pochodna $Du$ jako miara wektorowa} & \ref*{thm:BV-measure-29} & \pageref*{thm:BV-measure-29} \\
\hyperref[thm:RN-inverse-18]{Pochodna odwrotna} & \ref*{thm:RN-inverse-18} & \pageref*{thm:RN-inverse-18} \\
\hyperref[thm:higher-differentiation-11]{Pochodne wyższego rzędu} & \ref*{thm:higher-differentiation-11} & \pageref*{thm:higher-differentiation-11} \\
\hyperref[thm:weak-subseq-15]{Podciąg słabo zbieżny: wariant separowalny} & \ref*{thm:weak-subseq-15} & \pageref*{thm:weak-subseq-15} \\
\hyperref[thm:measure-subsequence]{Podciąg zbieżny prawie wszędzie} & \ref*{thm:measure-subsequence} & \pageref*{thm:measure-subsequence} \\
\hyperref[thm:conv-implications-9]{Podstawowe implikacje} & \ref*{thm:conv-implications-9} & \pageref*{thm:conv-implications-9} \\
\hyperref[thm:FTC-lebesgue-25]{Podstawowe twierdzenie rachunku całkowego dla całki Lebesgue'a} & \ref*{thm:FTC-lebesgue-25} & \pageref*{thm:FTC-lebesgue-25} \\
\hyperref[thm:basic-measure]{Podstawowe własności miary} & \ref*{thm:basic-measure} & \pageref*{thm:basic-measure} \\
\hyperref[thm:cantor-properties]{Podstawowe własności zbioru Cantora} & \ref*{thm:cantor-properties} & \pageref*{thm:cantor-properties} \\
\hyperref[thm:vitali-cover-24]{Pokryciowe twierdzenie Vitaliego} & \ref*{thm:vitali-cover-24} & \pageref*{thm:vitali-cover-24} \\
\hyperref[thm:portmanteau-22]{Portmanteau} & \ref*{thm:portmanteau-22} & \pageref*{thm:portmanteau-22} \\
\hyperref[thm:haar-modular-23]{Prawa miara Haara i modularność} & \ref*{thm:haar-modular-23} & \pageref*{thm:haar-modular-23} \\
\hyperref[thm:ae-smooth-approx-10]{Prawie wszędzie gładka aproksymacja funkcji mierzalnej} & \ref*{thm:ae-smooth-approx-10} & \pageref*{thm:ae-smooth-approx-10} \\
\hyperref[thm:ae-to-measure-finite]{Prawie wszędzie implikuje zbieżność według miary przy mierze skończonej} & \ref*{thm:ae-to-measure-finite} & \pageref*{thm:ae-to-measure-finite} \\
\hyperref[thm:volume-premeasure]{Premiara na ograniczonym prostopadłościanie} & \ref*{thm:volume-premeasure} & \pageref*{thm:volume-premeasure} \\
\hyperref[thm:prokhorov-22]{Prochorow na przestrzeni polskiej} & \ref*{thm:prokhorov-22} & \pageref*{thm:prokhorov-22} \\
\hyperref[thm:prokhorov-forward-22]{Prochorow: kierunek ciasność $\Rightarrow$ zwartość} & \ref*{thm:prokhorov-forward-22} & \pageref*{thm:prokhorov-forward-22} \\
\hyperref[thm:borel-product]{Produkt borelowski przestrzeni drugoprzeliczalnych} & \ref*{thm:borel-product} & \pageref*{thm:borel-product} \\
\hyperref[thm:lebesgue-products-12]{Produkt miar Lebesgue'a} & \ref*{thm:lebesgue-products-12} & \pageref*{thm:lebesgue-products-12} \\
\hyperref[thm:compact-extension-10]{Przedłużenie z kompaktu w $\mathbb R^d$} & \ref*{thm:compact-extension-10} & \pageref*{thm:compact-extension-10} \\
\hyperref[thm:intersection-sigma]{Przekroje $\sigma$-algebr} & \ref*{thm:intersection-sigma} & \pageref*{thm:intersection-sigma} \\
\hyperref[thm:borel-generators-R]{Przeliczalne generatory na $\mathbb R$} & \ref*{thm:borel-generators-R} & \pageref*{thm:borel-generators-R} \\
\hyperref[thm:standard-countably-generated]{Przeliczalne generowanie standardowych przestrzeni} & \ref*{thm:standard-countably-generated} & \pageref*{thm:standard-countably-generated} \\
\hyperref[thm:hilbert-rnp-30]{Przestrzenie Hilberta mają RNP} & \ref*{thm:hilbert-rnp-30} & \pageref*{thm:hilbert-rnp-30} \\
\hyperref[thm:CE-orthogonal-18]{Przypadek $L^2$: rzut ortogonalny} & \ref*{thm:CE-orthogonal-18} & \pageref*{thm:CE-orthogonal-18} \\
\hyperref[thm:simple-properties]{Rachunek dla funkcji prostych} & \ref*{thm:simple-properties} & \pageref*{thm:simple-properties} \\
\hyperref[thm:rademacher-28]{Rademacher} & \ref*{thm:rademacher-28} & \pageref*{thm:rademacher-28} \\
\hyperref[thm:RN-finite-18]{Radon--Nikodym dla miar skończonych} & \ref*{thm:RN-finite-18} & \pageref*{thm:RN-finite-18} \\
\hyperref[thm:RN-complex-18]{Radon--Nikodym dla miar zespolonych} & \ref*{thm:RN-complex-18} & \pageref*{thm:RN-complex-18} \\
\hyperref[thm:reflexive-lp-15]{Refleksywność $L^p$} & \ref*{thm:reflexive-lp-15} & \pageref*{thm:reflexive-lp-15} \\
\hyperref[thm:hausdorff-gdelta-26]{Regularizacja przez zbiory $G_\delta$} & \ref*{thm:hausdorff-gdelta-26} & \pageref*{thm:hausdorff-gdelta-26} \\
\hyperref[thm:RN-chain-18]{Reguła łańcuchowa} & \ref*{thm:RN-chain-18} & \pageref*{thm:RN-chain-18} \\
\hyperref[thm:riesz-positive-21]{Reprezentacja dodatnich funkcjonałów} & \ref*{thm:riesz-positive-21} & \pageref*{thm:riesz-positive-21} \\
\hyperref[thm:compact-riesz-21]{Reprezentacja na zwartej przestrzeni metrycznej} & \ref*{thm:compact-riesz-21} & \pageref*{thm:compact-riesz-21} \\
\hyperref[thm:riesz-complex-21]{Reprezentacja zespolona i norma wariacji} & \ref*{thm:riesz-complex-21} & \pageref*{thm:riesz-complex-21} \\
\hyperref[thm:riesz-fischer-14]{Riesz--Fischer: zupełność $L^p$} & \ref*{thm:riesz-fischer-14} & \pageref*{thm:riesz-fischer-14} \\
\hyperref[thm:rnp-lipschitz-equiv-30]{RNP a różniczkowalność krzywych Lipschitzowskich} & \ref*{thm:rnp-lipschitz-equiv-30} & \pageref*{thm:rnp-lipschitz-equiv-30} \\
\hyperref[thm:polar-complex-measure-18]{Rozkład biegunowy miary zespolonej} & \ref*{thm:polar-complex-measure-18} & \pageref*{thm:polar-complex-measure-18} \\
\hyperref[thm:hahn-16]{Rozkład Hahna} & \ref*{thm:hahn-16} & \pageref*{thm:hahn-16} \\
\hyperref[thm:jordan-16]{Rozkład Jordana} & \ref*{thm:jordan-16} & \pageref*{thm:jordan-16} \\
\hyperref[thm:lebesgue-decomp-finite-17]{Rozkład Lebesgue'a dla miar skończonych} & \ref*{thm:lebesgue-decomp-finite-17} & \pageref*{thm:lebesgue-decomp-finite-17} \\
\hyperref[thm:lebesgue-decomp-sigfin-17]{Rozkład Lebesgue'a w przypadku $\sigma$-skończonym} & \ref*{thm:lebesgue-decomp-sigfin-17} & \pageref*{thm:lebesgue-decomp-sigfin-17} \\
\hyperref[thm:leb-manifold-27]{Rozkład na część objętościową i osobliwą} & \ref*{thm:leb-manifold-27} & \pageref*{thm:leb-manifold-27} \\
\hyperref[thm:extension-caratheodory]{Rozszerzenie Carathéodory'ego} & \ref*{thm:extension-caratheodory} & \pageref*{thm:extension-caratheodory} \\
\hyperref[thm:measure-differentiation-25]{Różniczkowanie miary względem Lebesgue'a} & \ref*{thm:measure-differentiation-25} & \pageref*{thm:measure-differentiation-25} \\
\hyperref[thm:differentiating-integrals]{Różniczkowanie pod znakiem całki} & \ref*{thm:differentiating-integrals} & \pageref*{thm:differentiating-integrals} \\
\hyperref[thm:hausdorff-borel-26]{Samodzielna konstrukcja miary Hausdorffa} & \ref*{thm:hausdorff-borel-26} & \pageref*{thm:hausdorff-borel-26} \\
\hyperref[thm:separable-lp-14]{Separowalność $L^p$} & \ref*{thm:separable-lp-14} & \pageref*{thm:separable-lp-14} \\
\hyperref[thm:hausdorff-lipschitz-26]{Skalowanie i odwzorowania Lipschitzowskie} & \ref*{thm:hausdorff-lipschitz-26} & \pageref*{thm:hausdorff-lipschitz-26} \\
\hyperref[thm:besicovitch-finite-24]{Skończona postać lematu Besicovitcha} & \ref*{thm:besicovitch-finite-24} & \pageref*{thm:besicovitch-finite-24} \\
\hyperref[thm:measure-operations]{Sumy i mnożenie przez stałą} & \ref*{thm:measure-operations} & \pageref*{thm:measure-operations} \\
\hyperref[thm:sum-nonnegative-integrals]{Szeregi nieujemnych funkcji mierzalnych} & \ref*{thm:sum-nonnegative-integrals} & \pageref*{thm:sum-nonnegative-integrals} \\
\hyperref[thm:extended-tests]{Testy mierzalności na półprostych} & \ref*{thm:extended-tests} & \pageref*{thm:extended-tests} \\
\hyperref[thm:complex-test]{Test na częściach rzeczywistej i urojonej} & \ref*{thm:complex-test} & \pageref*{thm:complex-test} \\
\hyperref[thm:tonelli-12]{Tonelli} & \ref*{thm:tonelli-12} & \pageref*{thm:tonelli-12} \\
\hyperref[thm:completed-product-fubini-12]{Tonelli i Fubini dla uzupełnionej miary produktowej} & \ref*{thm:completed-product-fubini-12} & \pageref*{thm:completed-product-fubini-12} \\
\hyperref[thm:section-identity-12]{Tożsamość przekrojowa przy $\sigma$-skończoności} & \ref*{thm:section-identity-12} & \pageref*{thm:section-identity-12} \\
\hyperref[thm:density-transform-13]{Transformacja gęstości} & \ref*{thm:density-transform-13} & \pageref*{thm:density-transform-13} \\
\hyperref[thm:invariance-lebesgue]{Translacje i jednolite skalowanie} & \ref*{thm:invariance-lebesgue} & \pageref*{thm:invariance-lebesgue} \\
\hyperref[thm:caratheodory]{Twierdzenie Carathéodory'ego} & \ref*{thm:caratheodory} & \pageref*{thm:caratheodory} \\
\hyperref[thm:gauss-green-BV-29]{Twierdzenie Gaussa--Greena w postaci miarowej} & \ref*{thm:gauss-green-BV-29} & \pageref*{thm:gauss-green-BV-29} \\
\hyperref[thm:monotone-class]{Twierdzenie o klasie monotonicznej} & \ref*{thm:monotone-class} & \pageref*{thm:monotone-class} \\
\hyperref[thm:pi-lambda]{Twierdzenie $\pi$--$\lambda$ Dynkina} & \ref*{thm:pi-lambda} & \pageref*{thm:pi-lambda} \\
\hyperref[thm:besicovitch-24]{Twierdzenie pokryciowe Besicovitcha -- pełny dowód} & \ref*{thm:besicovitch-24} & \pageref*{thm:besicovitch-24} \\
\hyperref[thm:RN-sigfin-18]{Twierdzenie Radona--Nikodyma} & \ref*{thm:RN-sigfin-18} & \pageref*{thm:RN-sigfin-18} \\
\hyperref[thm:de-giorgi-structure-29]{Twierdzenie strukturalne De Giorgiego} & \ref*{thm:de-giorgi-structure-29} & \pageref*{thm:de-giorgi-structure-29} \\
\hyperref[thm:completion]{Uzupełnienie przestrzeni z miarą} & \ref*{thm:completion} & \pageref*{thm:completion} \\
\hyperref[thm:caratheodory-completion]{Uzupełnienie rozszerzenia} & \ref*{thm:caratheodory-completion} & \pageref*{thm:caratheodory-completion} \\
\hyperref[thm:vector-variation-measure-30]{Wariacja jest miarą dodatnią} & \ref*{thm:vector-variation-measure-30} & \pageref*{thm:vector-variation-measure-30} \\
\hyperref[thm:density-variation-16]{Wariacja miary określonej całką} & \ref*{thm:density-variation-16} & \pageref*{thm:density-variation-16} \\
\hyperref[thm:variation-density-30]{Wariacja miary o gęstości Bochnera} & \ref*{thm:variation-density-30} & \pageref*{thm:variation-density-30} \\
\hyperref[thm:complex-variation-16]{Wariacja miary zespolonej} & \ref*{thm:complex-variation-16} & \pageref*{thm:complex-variation-16} \\
\hyperref[thm:conditional-jensen-19]{Warunkowa nierówność Jensena} & \ref*{thm:conditional-jensen-19} & \pageref*{thm:conditional-jensen-19} \\
\hyperref[thm:lp-embedding-finite-14]{Włożenia przy skończonej mierze} & \ref*{thm:lp-embedding-finite-14} & \pageref*{thm:lp-embedding-finite-14} \\
\hyperref[thm:cantor-dimension-26]{Wymiar klasycznego zbioru Cantora} & \ref*{thm:cantor-dimension-26} & \pageref*{thm:cantor-dimension-26} \\
\hyperref[thm:euclidean-dimension-26]{Wymiar przestrzeni euklidesowej} & \ref*{thm:euclidean-dimension-26} & \pageref*{thm:euclidean-dimension-26} \\
\hyperref[thm:lipschitz-area-coarea-28]{Wzory area i coarea dla odwzorowań Lipschitzowskich} & \ref*{thm:lipschitz-area-coarea-28} & \pageref*{thm:lipschitz-area-coarea-28} \\
\hyperref[thm:coarea-submersion-28]{Wzór koobszaru dla submersji} & \ref*{thm:coarea-submersion-28} & \pageref*{thm:coarea-submersion-28} \\
\hyperref[thm:layer-cake-8]{Wzór warstwowy na całkę nieujemną} & \ref*{thm:layer-cake-8} & \pageref*{thm:layer-cake-8} \\
\hyperref[thm:immersion-area-28]{Wzór z krotnością dla immersji $C^1$} & \ref*{thm:immersion-area-28} & \pageref*{thm:immersion-area-28} \\
\hyperref[thm:hausdorff-line-26]{W wymiarze jeden otrzymujemy miarę Lebesgue'a} & \ref*{thm:hausdorff-line-26} & \pageref*{thm:hausdorff-line-26} \\
\hyperref[thm:change-variables-13]{Zamiana zmiennych dla dyfeomorfizmu $C^1$} & \ref*{thm:change-variables-13} & \pageref*{thm:change-variables-13} \\
\hyperref[thm:manifold-change-27]{Zamiana zmiennych na rozmaitościach} & \ref*{thm:manifold-change-27} & \pageref*{thm:manifold-change-27} \\
\hyperref[thm:kl-data-processing-19]{Zasada przetwarzania danych dla entropii} & \ref*{thm:kl-data-processing-19} & \pageref*{thm:kl-data-processing-19} \\
\hyperref[thm:mass-distribution-26]{Zasada rozkładu masy} & \ref*{thm:mass-distribution-26} & \pageref*{thm:mass-distribution-26} \\
\hyperref[thm:uniform-parameter-11]{Zbieżność jednostajna po scałkowaniu} & \ref*{thm:uniform-parameter-11} & \pageref*{thm:uniform-parameter-11} \\
\hyperref[thm:mct-ae]{Zbieżność monotoniczna prawie wszędzie} & \ref*{thm:mct-ae} & \pageref*{thm:mct-ae} \\
\hyperref[thm:vitali-l1]{Zbieżność Vitaliego na przestrzeni o mierze skończonej} & \ref*{thm:vitali-l1} & \pageref*{thm:vitali-l1} \\
\hyperref[thm:lp-to-measure-14]{Zbieżność w $L^p$ a zbieżność według miary} & \ref*{thm:lp-to-measure-14} & \pageref*{thm:lp-to-measure-14} \\
\hyperref[thm:tv-implies-narrow-22]{Zbieżność w wariacji implikuje wąską} & \ref*{thm:tv-implies-narrow-22} & \pageref*{thm:tv-implies-narrow-22} \\
\hyperref[thm:borel-trace]{Zbiory borelowskie podprzestrzeni} & \ref*{thm:borel-trace} & \pageref*{thm:borel-trace} \\
\hyperref[thm:riemann-to-lebesgue]{Zgodność całki Riemanna i Lebesgue'a} & \ref*{thm:riemann-to-lebesgue} & \pageref*{thm:riemann-to-lebesgue} \\
\hyperref[thm:composition-measurable]{Złożenia i funkcje ciągłe} & \ref*{thm:composition-measurable} & \pageref*{thm:composition-measurable} \\
\hyperref[thm:BV-banach-29]{Zupełność przestrzeni $BV$} & \ref*{thm:BV-banach-29} & \pageref*{thm:BV-banach-29} \\
\hyperref[thm:complete-measures-16]{Zupełność przestrzeni miar} & \ref*{thm:complete-measures-16} & \pageref*{thm:complete-measures-16} \\
\hyperref[thm:cauchy-measure-complete-9]{Zupełność zbieżności według miary} & \ref*{thm:cauchy-measure-complete-9} & \pageref*{thm:cauchy-measure-complete-9} \\
\hyperref[thm:perimeter-compact-29]{Zwartość rodzin zbiorów o ograniczonym obwodzie} & \ref*{thm:perimeter-compact-29} & \pageref*{thm:perimeter-compact-29} \\
\end{longtable}
\endgroup

\section*{Lematy}
\begingroup\small
\setlength\LTleft{0pt}\setlength\LTright{0pt}
\begin{longtable}{@{}p{0.72\linewidth}p{0.135\linewidth}p{0.08\linewidth}@{}}
\textbf{Nazwa} & \textbf{Nr} & \textbf{Str.}\\\hline
\endfirsthead
\multicolumn{3}{@{}l}{\emph{Lematy -- ciąg dalszy}}\\[2pt]
\textbf{Nazwa} & \textbf{Nr} & \textbf{Str.}\\\hline
\endhead
\hyperref[lem:hausdorff-metricity-26]{Addytywność na zbiorach oddalonych} & \ref*{lem:hausdorff-metricity-26} & \pageref*{lem:hausdorff-metricity-26} \\
\hyperref[lem:outer-approx-algebra]{Aproksymacja od góry} & \ref*{lem:outer-approx-algebra} & \pageref*{lem:outer-approx-algebra} \\
\hyperref[lem:vector-variation-ac-30]{Bezwzględna ciągłość wariacji} & \ref*{lem:vector-variation-ac-30} & \pageref*{lem:vector-variation-ac-30} \\
\hyperref[lem:metric-outer-borel-21]{Borelowska mierzalność} & \ref*{lem:metric-outer-borel-21} & \pageref*{lem:metric-outer-borel-21} \\
\hyperref[lem:current-boundary-square-29]{Brzeg brzegu znika} & \ref*{lem:current-boundary-square-29} & \pageref*{lem:current-boundary-square-29} \\
\hyperref[lem:essential-boundary-null-29]{Brzeg miarowy jest objętościowo mały} & \ref*{lem:essential-boundary-null-29} & \pageref*{lem:essential-boundary-null-29} \\
\hyperref[lem:FTC-Banach-30]{Całka po krzywej z pochodną Banachowską} & \ref*{lem:FTC-Banach-30} & \pageref*{lem:FTC-Banach-30} \\
\hyperref[lem:integration-density-18]{Całkowanie względem miary o gęstości} & \ref*{lem:integration-density-18} & \pageref*{lem:integration-density-18} \\
\hyperref[lem:positive-content-26]{Dodatnia zawartość zbioru zwartego} & \ref*{lem:positive-content-26} & \pageref*{lem:positive-content-26} \\
\hyperref[lem:midpoint-lp-15]{Dwie nierówności dla punktu środkowego} & \ref*{lem:midpoint-lp-15} & \pageref*{lem:midpoint-lp-15} \\
\hyperref[lem:nonatomic-divisibility-30]{Dzielenie miary bezatomowej} & \ref*{lem:nonatomic-divisibility-30} & \pageref*{lem:nonatomic-divisibility-30} \\
\hyperref[lem:singularity-17]{Elementarne własności osobliwości} & \ref*{lem:singularity-17} & \pageref*{lem:singularity-17} \\
\hyperref[lem:haar-fubini-compact-23]{Fubini dla ciągłych jąder o zwartym nośniku} & \ref*{lem:haar-fubini-compact-23} & \pageref*{lem:haar-fubini-compact-23} \\
\hyperref[lem:cutoff-radon-20]{Funkcja odcinająca} & \ref*{lem:cutoff-radon-20} & \pageref*{lem:cutoff-radon-20} \\
\hyperref[lem:density-cc-c0-21]{Gęstość $C_c(X)$ w $C_0(X)$} & \ref*{lem:density-cc-c0-21} & \pageref*{lem:density-cc-c0-21} \\
\hyperref[lem:dg-density-identification]{Gęstość miary na wykresach prostowalnych} & \ref*{lem:dg-density-identification} & \pageref*{lem:dg-density-identification} \\
\hyperref[lem:smooth-cutoff-10]{Gładka funkcja odcinająca} & \ref*{lem:smooth-cutoff-10} & \pageref*{lem:smooth-cutoff-10} \\
\hyperref[lem:canonical-lpq-15]{Izometryczne zanurzenie do dualnej} & \ref*{lem:canonical-lpq-15} & \pageref*{lem:canonical-lpq-15} \\
\hyperref[lem:haar-ratio-23]{Jednolite oszacowania ilorazów} & \ref*{lem:haar-ratio-23} & \pageref*{lem:haar-ratio-23} \\
\hyperref[lem:weak-sequence-tight-22]{Jednostajna ciasność ciągu wąsko zbieżnego} & \ref*{lem:weak-sequence-tight-22} & \pageref*{lem:weak-sequence-tight-22} \\
\hyperref[lem:haar-uniform-23]{Jednostajna ciągłość translacji funkcji} & \ref*{lem:haar-uniform-23} & \pageref*{lem:haar-uniform-23} \\
\hyperref[lem:rademacher-line-28]{Jednowymiarowa różniczkowalność} & \ref*{lem:rademacher-line-28} & \pageref*{lem:rademacher-line-28} \\
\hyperref[lem:uniform-absolute-continuity]{Kontrola całek na małych zbiorach} & \ref*{lem:uniform-absolute-continuity} & \pageref*{lem:uniform-absolute-continuity} \\
\hyperref[lem:hausdorff-caratheodory-selfcontained-26]{Kryterium Carathéodory'ego, z dowodem} & \ref*{lem:hausdorff-caratheodory-selfcontained-26} & \pageref*{lem:hausdorff-caratheodory-selfcontained-26} \\
\hyperref[lem:dynkin-intersection]{Lemat z rozdziału 1} & \ref*{lem:dynkin-intersection} & \pageref*{lem:dynkin-intersection} \\
\hyperref[lem:lipschitz-approx-compact-10]{Lipschitzowska aproksymacja na zbiorze zwartym} & \ref*{lem:lipschitz-approx-compact-10} & \pageref*{lem:lipschitz-approx-compact-10} \\
\hyperref[lem:lusin-simple-10]{Łuzin dla funkcji prostych} & \ref*{lem:lusin-simple-10} & \pageref*{lem:lusin-simple-10} \\
\hyperref[lem:haar-neigh-23]{Małe otoczenia, zwarte rozdrobnienia i pokrycia} & \ref*{lem:haar-neigh-23} & \pageref*{lem:haar-neigh-23} \\
\hyperref[lem:graph-hausdorff-28]{Małe zniekształcenie wykresu} & \ref*{lem:graph-hausdorff-28} & \pageref*{lem:graph-hausdorff-28} \\
\hyperref[lem:hausdorff-metric-open-26]{Metryczna miara zewnętrzna a zbiory domknięte} & \ref*{lem:hausdorff-metric-open-26} & \pageref*{lem:hausdorff-metric-open-26} \\
\hyperref[lem:BV-level-meas-auditVI]{Mierzalność obwodów poziomic} & \ref*{lem:BV-level-meas-auditVI} & \pageref*{lem:BV-level-meas-auditVI} \\
\hyperref[lem:sections-measurable-12]{Mierzalność wszystkich przekrojów} & \ref*{lem:sections-measurable-12} & \pageref*{lem:sections-measurable-12} \\
\hyperref[lem:variational-lp-15]{Minimum energii} & \ref*{lem:variational-lp-15} & \pageref*{lem:variational-lp-15} \\
\hyperref[lem:dg-relative-isoperimetric]{Nierówności Sobolewa i względna nierówność izoperymetryczna} & \ref*{lem:dg-relative-isoperimetric} & \pageref*{lem:dg-relative-isoperimetric} \\
\hyperref[lem:rect-cover-ineq]{Nierówność objętościowa dla pokryć} & \ref*{lem:rect-cover-ineq} & \pageref*{lem:rect-cover-ineq} \\
\hyperref[lem:simple-refinement]{Niezależność od reprezentacji} & \ref*{lem:simple-refinement} & \pageref*{lem:simple-refinement} \\
\hyperref[lem:critical-area-null-28]{Obraz zbioru o niedostatecznej randze} & \ref*{lem:critical-area-null-28} & \pageref*{lem:critical-area-null-28} \\
\hyperref[lem:BV-smooth-levels]{Obwody poziomic funkcji gładkiej} & \ref*{lem:BV-smooth-levels} & \pageref*{lem:BV-smooth-levels} \\
\hyperref[lem:lipschitz-null-13]{Odwzorowania Lipschitzowskie zachowują zbiory zerowe} & \ref*{lem:lipschitz-null-13} & \pageref*{lem:lipschitz-null-13} \\
\hyperref[lem:haar-content-to-measure-23]{Od zawartości na kompaktach do miary borelowskiej} & \ref*{lem:haar-content-to-measure-23} & \pageref*{lem:haar-content-to-measure-23} \\
\hyperref[lem:bounded-signed-16]{Ograniczoność zakresu} & \ref*{lem:bounded-signed-16} & \pageref*{lem:bounded-signed-16} \\
\hyperref[lem:dg-density]{Oszacowania gęstości na brzegu zredukowanym} & \ref*{lem:dg-density} & \pageref*{lem:dg-density} \\
\hyperref[lem:small-cube-13]{Oszacowanie obrazu małego sześcianu} & \ref*{lem:small-cube-13} & \pageref*{lem:small-cube-13} \\
\hyperref[lem:BV-translation-29]{Oszacowanie przesunięcia} & \ref*{lem:BV-translation-29} & \pageref*{lem:BV-translation-29} \\
\hyperref[lem:perimeter-complement-29]{Pochodna i dopełnienie} & \ref*{lem:perimeter-complement-29} & \pageref*{lem:perimeter-complement-29} \\
\hyperref[lem:premeasure-cover]{Podaddytywność premiary na pokryciach} & \ref*{lem:premeasure-cover} & \pageref*{lem:premeasure-cover} \\
\hyperref[lem:rectifiable-union-29]{Podstawowe własności prostowalności} & \ref*{lem:rectifiable-union-29} & \pageref*{lem:rectifiable-union-29} \\
\hyperref[lem:dg-blowup]{Półprzestrzeń jako jedyna granica po przeskalowaniu} & \ref*{lem:dg-blowup} & \pageref*{lem:dg-blowup} \\
\hyperref[lem:dg-rectifiability]{Prostowalność brzegu zredukowanego} & \ref*{lem:dg-rectifiability} & \pageref*{lem:dg-rectifiability} \\
\hyperref[lem:hausdorff-threshold-26]{Próg zerowania} & \ref*{lem:hausdorff-threshold-26} & \pageref*{lem:hausdorff-threshold-26} \\
\hyperref[lem:dg-ball-cut]{Przecięcie zbioru skończonego obwodu z kulą} & \ref*{lem:dg-ball-cut} & \pageref*{lem:dg-ball-cut} \\
\hyperref[lem:whitney-c1-28]{Przedłużenie zgodnych danych pierwszego rzędu} & \ref*{lem:whitney-c1-28} & \pageref*{lem:whitney-c1-28} \\
\hyperref[lem:finite-rectangle-sections-12]{Przekroje na prostokącie skończonej miary} & \ref*{lem:finite-rectangle-sections-12} & \pageref*{lem:finite-rectangle-sections-12} \\
\hyperref[lem:lipschitz-extend-22]{Przybliżanie lipschitzowskie i przedłużanie} & \ref*{lem:lipschitz-extend-22} & \pageref*{lem:lipschitz-extend-22} \\
\hyperref[lem:continuous-dense-25]{Przybliżanie lokalne funkcją ciągłą} & \ref*{lem:continuous-dense-25} & \pageref*{lem:continuous-dense-25} \\
\hyperref[lem:rnp-partition-30]{Przybliżanie miary przez skończone podziały} & \ref*{lem:rnp-partition-30} & \pageref*{lem:rnp-partition-30} \\
\hyperref[lem:compact-inside-10]{Przybliżenie mierzalnego zbioru przez zbiór zwarty} & \ref*{lem:compact-inside-10} & \pageref*{lem:compact-inside-10} \\
\hyperref[lem:radon-sigmafinite-inner-20]{Regularność na zbiorach $\sigma$-skończonych} & \ref*{lem:radon-sigmafinite-inner-20} & \pageref*{lem:radon-sigmafinite-inner-20} \\
\hyperref[lem:bochner-differentiation-30]{Różniczkowanie całki Bochnera} & \ref*{lem:bochner-differentiation-30} & \pageref*{lem:bochner-differentiation-30} \\
\hyperref[lem:besicovitch-angle-24]{Skończoność konfiguracji Besicovitcha} & \ref*{lem:besicovitch-angle-24} & \pageref*{lem:besicovitch-angle-24} \\
\hyperref[lem:rect-grid]{Skończony podział na siatkę} & \ref*{lem:rect-grid} & \pageref*{lem:rect-grid} \\
\hyperref[lem:rademacher-direction-28]{Słabe i kierunkowe pochodne funkcji Lipschitzowskiej} & \ref*{lem:rademacher-direction-28} & \pageref*{lem:rademacher-direction-28} \\
\hyperref[lem:BV-strict-approx]{Ścisła aproksymacja gładka w $BV$} & \ref*{lem:BV-strict-approx} & \pageref*{lem:BV-strict-approx} \\
\hyperref[lem:trace-generation]{Ślad i generowanie} & \ref*{lem:trace-generation} & \pageref*{lem:trace-generation} \\
\hyperref[lem:linear-coarea-28]{Tożsamość liniowo-algebraiczna koobszaru} & \ref*{lem:linear-coarea-28} & \pageref*{lem:linear-coarea-28} \\
\hyperref[lem:shear-13]{Translacje, permutacje i ścinanie} & \ref*{lem:shear-13} & \pageref*{lem:shear-13} \\
\hyperref[lem:BV-smooth-29]{Wariacja gładkiej funkcji} & \ref*{lem:BV-smooth-29} & \pageref*{lem:BV-smooth-29} \\
\hyperref[lem:hilbert-cube-22]{Włożenie separowalnej przestrzeni metrycznej w zwartą przestrzeń metryczną} & \ref*{lem:hilbert-cube-22} & \pageref*{lem:hilbert-cube-22} \\
\hyperref[lem:critical-coarea-null-28]{Włókna zwartego zbioru krytycznego} & \ref*{lem:critical-coarea-null-28} & \pageref*{lem:critical-coarea-null-28} \\
\hyperref[lem:conditional-pullout-19]{Wyciąganie czynnika mierzalnego} & \ref*{lem:conditional-pullout-19} & \pageref*{lem:conditional-pullout-19} \\
\hyperref[lem:exhaustion-21]{Wyczerpanie przez zbiory relatywnie zwarte} & \ref*{lem:exhaustion-21} & \pageref*{lem:exhaustion-21} \\
\hyperref[lem:mollify-cc-10]{Wygładzanie funkcji ciągłej o zwartym nośniku} & \ref*{lem:mollify-cc-10} & \pageref*{lem:mollify-cc-10} \\
\hyperref[lem:c1-general-area-28]{Wzór area dla dowolnego odwzorowania $C^1$} & \ref*{lem:c1-general-area-28} & \pageref*{lem:c1-general-area-28} \\
\hyperref[lem:young-numbers-14]{Young dla iloczynu dwóch liczb} & \ref*{lem:young-numbers-14} & \pageref*{lem:young-numbers-14} \\
\hyperref[lem:singular-density-25]{Zanik gęstości miary osobliwej} & \ref*{lem:singular-density-25} & \pageref*{lem:singular-density-25} \\
\hyperref[lem:haar-content-23]{Zawartość niezmiennicza na kompaktach} & \ref*{lem:haar-content-23} & \pageref*{lem:haar-content-23} \\
\hyperref[lem:simple-exhaustion]{Zbieżność całek prostej funkcji na zbiorach rosnących} & \ref*{lem:simple-exhaustion} & \pageref*{lem:simple-exhaustion} \\
\hyperref[lem:hyperplane-null]{Zbiory zawarte w hiperpłaszczyźnie} & \ref*{lem:hyperplane-null} & \pageref*{lem:hyperplane-null} \\
\hyperref[lem:null-fibres-lipschitz-28]{Zbiory zerowe i ich włókna} & \ref*{lem:null-fibres-lipschitz-28} & \pageref*{lem:null-fibres-lipschitz-28} \\
\hyperref[lem:hausdorff-outer-26]{Zewnętrzne własności kosztu pokryć} & \ref*{lem:hausdorff-outer-26} & \pageref*{lem:hausdorff-outer-26} \\
\hyperref[lem:restriction-integral]{Zgodność z ograniczeniem miary} & \ref*{lem:restriction-integral} & \pageref*{lem:restriction-integral} \\
\hyperref[lem:esssup-14]{Znaczenie kresu istotnego} & \ref*{lem:esssup-14} & \pageref*{lem:esssup-14} \\
\hyperref[lem:lusin-whitney-28]{Zwarte fragmenty klasy $C^1$} & \ref*{lem:lusin-whitney-28} & \pageref*{lem:lusin-whitney-28} \\
\hyperref[lem:haar-product-compact-23]{Zwartość dowolnego produktu przedziałów} & \ref*{lem:haar-product-compact-23} & \pageref*{lem:haar-product-compact-23} \\
\hyperref[lem:compact-prob-22]{Zwartość prawdopodobieństw na przestrzeni zwartej} & \ref*{lem:compact-prob-22} & \pageref*{lem:compact-prob-22} \\
\end{longtable}
\endgroup

\section*{Wnioski}
\begingroup\small
\setlength\LTleft{0pt}\setlength\LTright{0pt}
\begin{longtable}{@{}p{0.72\linewidth}p{0.135\linewidth}p{0.08\linewidth}@{}}
\textbf{Nazwa} & \textbf{Nr} & \textbf{Str.}\\\hline
\endfirsthead
\multicolumn{3}{@{}l}{\emph{Wnioski -- ciąg dalszy}}\\[2pt]
\textbf{Nazwa} & \textbf{Nr} & \textbf{Str.}\\\hline
\endhead
\hyperref[cor:smooth-approx-in-measure-10]{Aproksymacja w mierze funkcjami gładkimi} & \ref*{cor:smooth-approx-in-measure-10} & \pageref*{cor:smooth-approx-in-measure-10} \\
\hyperref[cor:besicovitch-measure-infinite-24]{Besicovitch dla dodatnich miar borelowskich} & \ref*{cor:besicovitch-measure-infinite-24} & \pageref*{cor:besicovitch-measure-infinite-24} \\
\hyperref[cor:borel-representative-10]{Borelowska reprezentacja funkcji Lebesgue'a} & \ref*{cor:borel-representative-10} & \pageref*{cor:borel-representative-10} \\
\hyperref[cor:coordinate-measurable]{Borelowskie odwzorowania wektorowe} & \ref*{cor:coordinate-measurable} & \pageref*{cor:coordinate-measurable} \\
\hyperref[cor:meas-simple-characterization]{Charakteryzacja mierzalności przez funkcje proste} & \ref*{cor:meas-simple-characterization} & \pageref*{cor:meas-simple-characterization} \\
\hyperref[cor:coercive-tight-22]{Ciasność z oszacowania momentów} & \ref*{cor:coercive-tight-22} & \pageref*{cor:coercive-tight-22} \\
\hyperref[cor:uniform-abs-cont-int-17]{Ciągłość całki względem zbioru całkowania} & \ref*{cor:uniform-abs-cont-int-17} & \pageref*{cor:uniform-abs-cont-int-17} \\
\hyperref[cor:l1-param-11]{Ciągłość w normie $L^1$} & \ref*{cor:l1-param-11} & \pageref*{cor:l1-param-11} \\
\hyperref[cor:haar-open-positive-23]{Dodatniość miary zbiorów otwartych} & \ref*{cor:haar-open-positive-23} & \pageref*{cor:haar-open-positive-23} \\
\hyperref[cor:holder-sharp-15]{Dokładne maksimum w nierówności Höldera} & \ref*{cor:holder-sharp-15} & \pageref*{cor:holder-sharp-15} \\
\hyperref[cor:lusin-finite-ae-10]{Funkcje skończone prawie wszędzie} & \ref*{cor:lusin-finite-ae-10} & \pageref*{cor:lusin-finite-ae-10} \\
\hyperref[cor:lebesgue-regularity-characterization]{Geometryczna charakterystyka mierzalności} & \ref*{cor:lebesgue-regularity-characterization} & \pageref*{cor:lebesgue-regularity-characterization} \\
\hyperref[cor:countable-lusin-10]{Jednoczesne twierdzenie Łuzina dla przeliczalnej rodziny} & \ref*{cor:countable-lusin-10} & \pageref*{cor:countable-lusin-10} \\
\hyperref[cor:measure-limit-unique]{Jednoznaczność granicy według miary} & \ref*{cor:measure-limit-unique} & \pageref*{cor:measure-limit-unique} \\
\hyperref[cor:vague-c0-bounded-22]{Kiedy zbieżność mglista jest słabą-*?} & \ref*{cor:vague-c0-bounded-22} & \pageref*{cor:vague-c0-bounded-22} \\
\hyperref[cor:NP-exact-alpha-19]{Konstrukcja testu optymalnego o zadanym poziomie} & \ref*{cor:NP-exact-alpha-19} & \pageref*{cor:NP-exact-alpha-19} \\
\hyperref[cor:simple-approx-integral]{Konstruktywne obliczanie całki} & \ref*{cor:simple-approx-integral} & \pageref*{cor:simple-approx-integral} \\
\hyperref[cor:conditional-lp-19]{Kontrakcja w $L^p$} & \ref*{cor:conditional-lp-19} & \pageref*{cor:conditional-lp-19} \\
\hyperref[cor:HL-level-24]{Kontrola zbioru, na którym średnie są duże} & \ref*{cor:HL-level-24} & \pageref*{cor:HL-level-24} \\
\hyperref[cor:tonelli-integrability-12]{Kryterium nieujemnej całkowalności} & \ref*{cor:tonelli-integrability-12} & \pageref*{cor:tonelli-integrability-12} \\
\hyperref[cor:practical-fubini-12]{Kryterium praktyczne} & \ref*{cor:practical-fubini-12} & \pageref*{cor:practical-fubini-12} \\
\hyperref[cor:every-subseq-9]{Kryterium za pomocą wszystkich podciągów} & \ref*{cor:every-subseq-9} & \pageref*{cor:every-subseq-9} \\
\hyperref[cor:nonborel-lebesgue]{Lebesgue'owskie zbiory nieborelowskie} & \ref*{cor:nonborel-lebesgue} & \pageref*{cor:nonborel-lebesgue} \\
\hyperref[cor:associative-product-12]{Łączność skończonych produktów} & \ref*{cor:associative-product-12} & \pageref*{cor:associative-product-12} \\
\hyperref[cor:riemann-hausdorff-28]{Miara riemannowska i Hausdorffa} & \ref*{cor:riemann-hausdorff-28} & \pageref*{cor:riemann-hausdorff-28} \\
\hyperref[cor:measures-separate-21]{Miary odróżniane przez funkcje ciągłe} & \ref*{cor:measures-separate-21} & \pageref*{cor:measures-separate-21} \\
\hyperref[cor:haar-adjoint-23]{Modularność grupy Liego} & \ref*{cor:haar-adjoint-23} & \pageref*{cor:haar-adjoint-23} \\
\hyperref[idx:06:wniosek:06]{Monotoniczne ciągi funkcji} & \ref*{idx:06:wniosek:06} & \pageref*{idx:06:wniosek:06} \\
\hyperref[cor:minimal-variation-16]{Najmniejsza miara dominująca} & \ref*{cor:minimal-variation-16} & \pageref*{cor:minimal-variation-16} \\
\hyperref[cor:integral-triangle]{Nierówność trójkąta i struktura $L^1$} & \ref*{cor:integral-triangle} & \pageref*{cor:integral-triangle} \\
\hyperref[cor:no-all-subsets-invariant]{Nie istnieje miara długości wszystkich podzbiorów} & \ref*{cor:no-all-subsets-invariant} & \pageref*{cor:no-all-subsets-invariant} \\
\hyperref[cor:submanifold-null-27]{Obrazy immersji niższego wymiaru} & \ref*{cor:submanifold-null-27} & \pageref*{cor:submanifold-null-27} \\
\hyperref[cor:reverse-fatou]{Odwrotny lemat Fatou przy dominacji} & \ref*{cor:reverse-fatou} & \pageref*{cor:reverse-fatou} \\
\hyperref[cor:5r-volume-24]{Oszacowanie miary sumy kul} & \ref*{cor:5r-volume-24} & \pageref*{cor:5r-volume-24} \\
\hyperref[cor:finite-besicovitch-measure-24]{Oszacowanie sumy miar w rodzinie skończonej} & \ref*{cor:finite-besicovitch-measure-24} & \pageref*{cor:finite-besicovitch-measure-24} \\
\hyperref[cor:RN-lebesgue-sum-18]{Pełny rozkład miary na część z gęstością i osobliwą} & \ref*{cor:RN-lebesgue-sum-18} & \pageref*{cor:RN-lebesgue-sum-18} \\
\hyperref[cor:riesz-order-21]{Porządek dodatnich funkcjonałów i miar} & \ref*{cor:riesz-order-21} & \pageref*{cor:riesz-order-21} \\
\hyperref[cor:null-sections-12]{Prawie wszystkie przekroje zbioru zerowego} & \ref*{cor:null-sections-12} & \pageref*{cor:null-sections-12} \\
\hyperref[cor:real-borel-tests]{Prawostronny test mierzalności} & \ref*{cor:real-borel-tests} & \pageref*{cor:real-borel-tests} \\
\hyperref[cor:null-union]{Przeliczalne sumy zbiorów miary zero} & \ref*{cor:null-union} & \pageref*{cor:null-union} \\
\hyperref[cor:countably-many-disjoint-atoms]{Przeliczalność rozłącznych atomów} & \ref*{cor:countably-many-disjoint-atoms} & \pageref*{cor:countably-many-disjoint-atoms} \\
\hyperref[idx:30:wniosek:15]{Przestrzenie Banacha i całki z krzywych} & \ref*{idx:30:wniosek:15} & \pageref*{idx:30:wniosek:15} \\
\hyperref[cor:lp-normed-14]{Przestrzeń normowana} & \ref*{cor:lp-normed-14} & \pageref*{cor:lp-normed-14} \\
\hyperref[cor:tv-data-processing-19]{Przetwarzanie danych nie zwiększa wariacji całkowitej} & \ref*{cor:tv-data-processing-19} & \pageref*{cor:tv-data-processing-19} \\
\hyperref[cor:indicator-smooth-10]{Przybliżanie funkcji charakterystycznej} & \ref*{cor:indicator-smooth-10} & \pageref*{cor:indicator-smooth-10} \\
\hyperref[cor:vitali-open-24]{Przybliżanie zbioru rozłącznymi kulami} & \ref*{cor:vitali-open-24} & \pageref*{cor:vitali-open-24} \\
\hyperref[cor:hausdorff-singletons-auditVI]{Punkty i zbiory przeliczalne} & \ref*{cor:hausdorff-singletons-auditVI} & \pageref*{cor:hausdorff-singletons-auditVI} \\
\hyperref[cor:integral-countable-partition]{Rozkład całki na rozłączne zbiory} & \ref*{cor:integral-countable-partition} & \pageref*{cor:integral-countable-partition} \\
\hyperref[cor:complex-lebesgue-decomp-17]{Rozkład miary zespolonej} & \ref*{cor:complex-lebesgue-decomp-17} & \pageref*{cor:complex-lebesgue-decomp-17} \\
\hyperref[cor:RN-linearity-18]{Różnica i suma miar z gęstościami} & \ref*{cor:RN-linearity-18} & \pageref*{cor:RN-linearity-18} \\
\hyperref[cor:coarea-level-28]{Różniczkowanie poziomic} & \ref*{cor:coarea-level-28} & \pageref*{cor:coarea-level-28} \\
\hyperref[cor:cubes-differentiation-25]{Średnie po kostkach} & \ref*{cor:cubes-differentiation-25} & \pageref*{cor:cubes-differentiation-25} \\
\hyperref[cor:bct]{Twierdzenie o zbieżności ograniczonej} & \ref*{cor:bct} & \pageref*{cor:bct} \\
\hyperref[idx:08:wniosek:20]{Uogólniona formuła rozkładu na poziomice} & \ref*{idx:08:wniosek:20} & \pageref*{idx:08:wniosek:20} \\
\hyperref[cor:generalized-holder-14]{Uogólniony Hölder} & \ref*{cor:generalized-holder-14} & \pageref*{cor:generalized-holder-14} \\
\hyperref[cor:diff-basis-25]{Wartość funkcji z otaczających ją średnich} & \ref*{cor:diff-basis-25} & \pageref*{cor:diff-basis-25} \\
\hyperref[idx:09:wniosek:14]{Wersja dla gęstości prawdopodobieństwa} & \ref*{idx:09:wniosek:14} & \pageref*{idx:09:wniosek:14} \\
\hyperref[cor:local-egorov-9]{Wersja lokalna dla przestrzeni $\sigma$-skończonej} & \ref*{cor:local-egorov-9} & \pageref*{cor:local-egorov-9} \\
\hyperref[cor:lusin-extension-10]{Wersja Łuzina przez funkcję ciągłą na całej przestrzeni} & \ref*{cor:lusin-extension-10} & \pageref*{cor:lusin-extension-10} \\
\hyperref[idx:06:wniosek:15]{Wersje mierzalne funkcji określonych p.w.} & \ref*{idx:06:wniosek:15} & \pageref*{idx:06:wniosek:15} \\
\hyperref[cor:graph-dimension-26]{Wykres funkcji Lipschitzowskiej} & \ref*{cor:graph-dimension-26} & \pageref*{cor:graph-dimension-26} \\
\hyperref[cor:hausdorff-union-26]{Wymiar sumy przeliczalnej} & \ref*{cor:hausdorff-union-26} & \pageref*{cor:hausdorff-union-26} \\
\hyperref[cor:BV-coarea-weighted-29]{Wzór koobszaru z wagą} & \ref*{cor:BV-coarea-weighted-29} & \pageref*{cor:BV-coarea-weighted-29} \\
\hyperref[cor:change-subdomain-13]{Zamiana na mierzalnym podzbiorze} & \ref*{cor:change-subdomain-13} & \pageref*{cor:change-subdomain-13} \\
\hyperref[cor:vague-narrow-22]{Zbieżność mglista bez ucieczki masy} & \ref*{cor:vague-narrow-22} & \pageref*{cor:vague-narrow-22} \\
\hyperref[cor:dct-l1]{Zbieżność zmajoryzowana w $L^1$} & \ref*{cor:dct-l1} & \pageref*{cor:dct-l1} \\
\hyperref[cor:measure-metric-complete-9]{Zupełność metryki zbieżności według miary} & \ref*{cor:measure-metric-complete-9} & \pageref*{cor:measure-metric-complete-9} \\
\end{longtable}
\endgroup


% Dżety dodane jako rozdział 31 — bez przenumerowania indeksu historycznego.
